\section{Appendix A. Mathematical Derivations}

\subsection{A.1 Recurrence Unrolling}

The decay recurrence:

\begin{equation}
S_t = \lambda_t S_{t-1} + \phi(k_t) v_t^\top
\end{equation}

With $S_0 = 0$:

\begin{equation}
S_t = \sum_{j=1}^{t} \left( \prod_{i=j+1}^{t} \lambda_i \right) \phi(k_j) v_j^\top
\end{equation}

\textbf{Decay coefficient:} $D_{t, j} = \exp(\sum_{i=j+1}^t \log \lambda_i)$.

\subsection{A.2 Chunked Scan Correctness}

\textbf{Claim.} The chunked scan produces the same output as the direct 
recurrence.

\textbf{Proof sketch.} By induction on chunks.

\textbf{Base case.} Chunk 0: no previous state. The intra-chunk computation 
matches the recurrence.

\textbf{Inductive step.} Assume chunk $c-1$ left state $(S_{c-1}, z_{c-1})$ 
matching the direct recurrence at the chunk boundary. Then:

\begin{itemize}
\item \textbf{Intra-chunk} computation uses the local cumsum, matching the recurrence
  within the chunk.
\item \textbf{Inter-chunk} contributions use the decayed state $D\_i\^{}\text{inter} =
  \exp(\Lambda\_i)$ applied to $S\_{c-1}$, which is exactly the decay from the
  boundary to position $i$.
\item \textbf{State update} produces $S_c$ via (4.16), which matches the recurrence
  at the chunk boundary.

\end{itemize}
$\square$

\subsection{A.3 ELU+1 Positivity}

\begin{equation}
\phi(x) = \text{elu}(x) + 1 > 0 \quad \forall x \in \mathbb{R}
\end{equation}

\textbf{Proof.} $\text{elu}(x) \in (-1, \infty)$, so $\phi(x) \in (0, \infty)$. 
$\square$

\subsection{A.4 Data-Dependent Decay}

\begin{equation}
\lambda_{h,t} = \sigma(W_g x_t + b_g + \ell_h)
\end{equation}

Initialization: $W_g = 0, b_g = 0$, $\ell_h \in [-0.5, 0.0]$.

\textbf{Head diversity.} $\lambda_{h} = \sigma(\ell_h)$ ranges:
\begin{itemize}
\item $\lambda_0 = \sigma(-0.5) \approx 0.38$
\item $\lambda_{H-1} = \sigma(0.0) \approx 0.50$

\end{itemize}
Effective horizons:
\begin{itemize}
\item Head 0: $1/(1-\lambda_0) \approx 1.6$
\item Head $H-1$: $1/(1-\lambda_{H-1}) \approx 2.0$

\end{itemize}
\subsection{A.5 Complexity}

For a single layer:

\begin{table}[H]
\centering
\small
\begin{tabular}{lll}
\toprule
\textbf{Component} & \textbf{Time} & \textbf{Memory} \\
\midrule
Projections & $\mathcal{O}(L d^2)$ & $\mathcal{O}(L d)$ \\
Cumsum & $\mathcal{O}(L)$ & $\mathcal{O}(L)$ \\
Intra-chunk & $\mathcal{O}(L \cdot C)$ & $\mathcal{O}(L \cdot C)$ \\
Inter-chunk & $\mathcal{O}(L \cdot d_h)$ & $\mathcal{O}(d_h^2)$ \\
State update & $\mathcal{O}(L \cdot d_h^2 / C)$ & $\mathcal{O}(d_h^2)$ \\
\bottomrule
\end{tabular}
\end{table}

\textbf{Total memory:} $\mathcal{O}(L d + L C + d_h^2)$.

For $C = 32$ and $d_h = 64$: state memory $= 16$ KB per layer.

\subsection{A.6 Bidirectional Combination}

\begin{equation}
O = \frac{1}{2} (O^\text{fwd} + O^\text{bwd})
\end{equation}

\textbf{Why 0.5?} Maintains output magnitude. Sum would double, average avoids it.

\subsection{A.7 Cross-Attention}

For encoder-decoder attention, K, V come from encoder states:

\begin{equation}
K = H_\text{enc} W_K, \quad V = H_\text{enc} W_V
\end{equation}

Q from decoder:

\begin{equation}
Q = H_\text{dec} W_Q
\end{equation}

Same chunked scan applies.

